\subsection{Emergent Fermions at the Endpoints of Strings -- Equivalence between Spin and Fermion Representation}\label{sec:Emergent-Fermions-at-the-Endpoints}

%\textcolor{red}{The first paragraph should be high level explaination first then detail descriptions later.}

\valert{The form of~\eqref{equ.density-matrix} implies the separability of the matter sector from the gauge sector via a pure spin representation. This observation also reflects the similarity between the density matrix represented using spins and the one using Majorana fermions. In this section, we demonstrate that the spin representation and the Majorana-fermion representation are identical, and they are just two different representations of the same algebra. Such an identity makes the density matrix reproduce the correct entanglement and further implies the block-diagonal structure of the density matrix.}

Starting from the Majorana representation, the ground state density matrix can also be represented as
\[
  \rho = 2^{N-1}\mathcal{P}_G\prod_{\anglelr{ij}}\lr{\frac{\mathds{1}+\hat{u}_{ij}}{2}}\rho_F\mathcal{P}_G
\]
where $\rho_F$ is the density matrix of the Kitaev ground state in the Majorana fermion representation. Details of the derivation are presented in Appendix \ref{appendix:Majorana-derivation}. The Majorana density matrix $\rho_F$ can be constructed from a starting point similar to \eqref{equ.rho-starting} that we have used in the beginning.
\begin{equation}\label{equ.rho-F}
    \rho^F = \frac{1}{2^{N/2}}\sum_{[\mathfrak{D}]}\mathfrak{C}_\mathfrak{D}\mathcal{M}_\mathfrak{D}
\end{equation}
where
\begin{equation}\label{equ.V-rep-def}
    \mathcal{M}_\mathfrak{D} \equiv \prod_{\anglelr{ij}\in\mathfrak{D}}ic_jc_i
\end{equation}
In each dimer $\anglelr{ij}$ of the above product, the site $i$ belongs to sublattice $A$ and $j$ to sublattice $B$. The denominator of the normalization coefficient is $2^{N/2}$ instead of $2^N$ because a normal fermion $f$ is defined by two $c$-Majorana fermions
%
%$c_1$ and $c_2$ as $f=\frac{1}{2}(c_1+ic_2)$ and $f^\dagger=\frac{1}{2}(c_1-ic_2)$. Therefore, the total number of physical degrees of freedom of a system composed of $N$ Majorana fermions is $2^{N/2}$ instead of $2^N$
%
. To manifest the similarity between the spin density matrix \eqref{equ.density-matrix}, we define a set of operators $\mathcal{S}_{\mathfrak{D}}$ for each operator that corresponds to an edge configuration
\begin{equation}\label{equ.S-def}
    \mathcal{S}_\mathfrak{D} \equiv \mathcal{P}^{(p)}\Sigma_\mathfrak{D}
\end{equation}
Given equation \eqref{equ.Sigma-representation} and the fact that $[\mathcal{P}_G,\Sigma_{\mathfrak{D}}]=0$, the set of operators $\lrg{\mathcal{S}_{\mathfrak{D}}}$ also forms a representation of the algebra of edge configurations $\anglelr{\lrg{\mathfrak{D}},\ominus}$, 
\begin{equation}\label{equ.U-rep-algebra}
    \mathcal{S}_{\mathfrak{D}}\mathcal{S}_{\mathfrak{D}'} = \varphi(\mathfrak{D},\mathfrak{D}')\mathcal{S}_{\mathfrak{D}\ominus\mathfrak{D}'}.
\end{equation}
With the string representation \eqref{equ.S-def} and \eqref{equ.U-rep-algebra}, The form of the density matrix in Eq.~\eqref{equ.density-matrix} becomes
\begin{equation}\label{equ.rho-U}
    \rho =  \frac{1}{2^{N-N_p}} \sum_{[\mathfrak{D}]}\mathfrak{C}_{\mathfrak{D}}\mathcal{S}_{\mathfrak{D}}
\end{equation}
which is of a form extremely similar to \eqref{equ.rho-F} if those $\mathcal{S}$-operators are replaced by $\mathcal{M}$-operators. This is reminiscent of different representations of the same algebra.
If we use the Majorana fermions to represent spin operators as $\sigma^\mu\to ib^\mu c$, then 
\begin{align*}
    \mathcal{S}_\mathfrak{D} &\to \mathcal{P}^{(p)} \prod_{\anglelr{ij}\in \mathfrak{D}}\lr{\hat{u}_{ij}ic_jc_i} \\
    &= \mathcal{P}^{(p)} \lrm{\prod_{\anglelr{ij}\in \mathfrak{D}}\hat{u}_{ij}} \mathcal{M}_{\mathfrak{D}}
\end{align*}
The first two products completely consist of $\hat{u}_{ij}$'s, which commute with each other as well as the $c$-Majorana fermion operators~\footnote{Note that the projector $\mathcal{P}^{(p)}$ becomes
\[\mathcal{P}^{(p)} \to \prod_{ p}\lr{\frac{\mathds{1} + \prod_{\anglelr{ij}\in p}\hat{u}_{ij}}{2}}\]
in the Majorana representation, where ``$\anglelr{ij}\in p$'' stands for \textit{those edges of the plaquette $p$}.}. Therefore, they will not contribute to the phase factor $g\lr{\mathfrak{D},\mathfrak{D}'}$ at all. So, we know that the anti-commutation relations among $c$-Majorana fermions are encoded in the phase factor $g\lr{\mathfrak{D},\mathfrak{D}'}$. In view of this, we find that the set of operators $\mathcal{M}_\mathfrak{D}$ defined by equation \eqref{equ.V-rep-def} also forms a projective representation of the edge configuration algebra and shares the same phase factor $g\lr{\mathfrak{D},\mathfrak{D}'}$ with the representation $\mathcal{S}$
\begin{equation}\label{equ.V-rep-algebra}
    \mathcal{M}_{\mathfrak{D}}\mathcal{M}_{\mathfrak{D}'} = \varphi(\mathfrak{D},\mathfrak{D}')\mathcal{M}_{\mathfrak{D}\ominus\mathfrak{D}'}
\end{equation}
This clearly illustrated a picture: the Majorana fermion representation is identical to the spin, or Pauli string, representation, and the Majorana fermions can be seen as emerging from the endpoints of those Pauli strings defined by \eqref{equ.string-operator-def}~\cite{Kitaev_2006, fermion-and-string, string-net, Photons-and-electrons-as-emergent-phenomena-string, PRD-string-wen-quantum-order}. First, we exploit the fact that the spin degree of freedom can be fractionalized into Majorana fermions, which allows us to solve the model analytically. Conversely, to represent everything in terms of spins, as our primary goal in this work, we must \textit{glue} the fractionalized degrees of freedom back via string operators defined in the spin language. 

\valert{
For the reduced density matrix, we can also define a set of $\mathcal{S}$-operators
\begin{equation}\label{equ.U-A-def}
    \mathcal{S}^{(A)}_\mathfrak{D} \equiv \mathcal{P}^{(p)}_A\Sigma_\mathfrak{D}
\end{equation}
supported on the subregion $A$. With these newly defined operators, the reduced density matrix becomes
\begin{equation}\label{equ.red-DM-new}
    \rho_A = \frac{1}{2^{N^A-N^A_p}}\sum_{[\mathfrak{D}]\subset A}\mathfrak{C}_{\mathfrak{D}}\mathcal{S}^{(A)}_\mathfrak{D}.
\end{equation}
The form of the reduced density matrix above implies that the physical picture of Majorana fermions emerging at the endpoints of string operators still holds in any simply-connected subregion $A$, and it can be verified that $\mathcal{S}^{(A)}$-operators also satisfy a relation
\begin{equation}\label{equ.U-A-rep-algebra}
    \mathcal{S}^{(A)}_{\mathfrak{D}}\mathcal{S}^{(A)}_{\mathfrak{D}'} = \varphi(\mathfrak{D},\mathfrak{D}')\mathcal{S}^{(A)}_{\mathfrak{D}\ominus\mathfrak{D}'},
\end{equation}
which is the same as \eqref{equ.U-rep-algebra} and \eqref{equ.V-rep-algebra}. The reduced density matrix of Majorana fermions in the region $A$ also possesses the same form 
\begin{equation}\label{equ.rho-F-reduced}
    \rho_A^{F} \equiv \frac{1}{2^{N^A/2}}\sum_{[\mathfrak{D}]\subset A}\mathfrak{C}_{\mathfrak{D}}\mathcal{M}_\mathfrak{D}
\end{equation}
There are two straightforward ways to obtain this conclusion. The first approach is to trace out the environment of the Majorana fermion density matrix \eqref{equ.rho-F}. Because every $\mathcal{M}_\mathfrak{D}$ is a $\mathbb{Z}_2$ operator and thus $\Tr \mathcal{M}_{\mathfrak{D}}=0$ if $\mathfrak{D}\not\sim\varnothing$, the remaining part therefore consists of terms with $[\mathfrak{D}]\subset A$, which constitutes the reduced density matrix of a simply-connected region $A$.  The second way is following the logic of \eqref{equ.rho-A-raw-guess}. The observables in the Majorana space with finite expectation values are exactly $\mathcal{M}_\mathfrak{D}$ operators. To ensure that the reduced density matrix $\rho^F_A$ preserves the correct expectation values for these operators, the form of $\rho^F_A$ must be of \eqref{equ.rho-F-reduced}. Here, we only consider regions with an even $N^A$ since the denominator $2^{N^A/2}$ becomes fractional if $N^A$ is an odd number, making the dimension of the Hilbert space ill-defined. In the latter sections, cases with an odd number of $N^A$ will be discussed. 
}


From the analytical form of the density matrices we found, we see the identity between the spin representation and the Majorana-fermion representation arises from the fact that every operator $\mathcal{S}_\mathfrak{D}$ depends only on the equivalence class $[\mathfrak{D}]$ the configuration $\mathfrak{D}$ belongs to. In other words, the following operators with different configurations, which share the same set of endpoints, are identical.
\[
\mathcal{S}\lr{\wordfig{0.5}{figures/ij_segment.pdf}} = \mathcal{S}\lr{\wordfig{0.75}{figures/ij_curve.pdf}} = \mathcal{S}\lr{\wordfig{0.8}{figures/ij_segment_1_loop.pdf}} = \cdots
\]
From the property of string operators $\Sigma_\mathfrak{D}$, we can directly prove the following equation
\begin{equation}\label{eq.S-rep-swap}
    \lsegment{i}{l}\lsegment{k}{ij} = -\lsegment{i}{j}\lsegment{k}{il},
\end{equation}
which can be represented pictorially as follows
\[
\mathcal{S}\lr{\wordfig{1.5}{figures/fermion-braiding_1.pdf}} = -\mathcal{S}\lr{\wordfig{1.5}{figures/fermion-braiding_2.pdf}}.
\]
The physical meaning of this equation is that the endpoints of strings obey fermion statistics~\cite{fermion-and-string}. The mathematical operation of this equation is equivalent to exchanging the positions of string-endpoints at sites $i$ and $k$. It turns out that an additional phase of $-1$ is accumulated after the swap, analogous to what happens when two fermions are exchanged. In fact, if we identify the operator $\lsegment{j}{i}$ of a single string $\overleftarrow{ji}$ pointing from $i$ to $j$ with the hopping operator $ic_jc_i$, which is actually the Majorana representation of the string $\mathcal{M}_{\overleftarrow{ji}}\propto ic_jc_i$, of Majorana fermions hopping from $i$ to $j$, then the above equation is equivalent to the following Majorana representation
\[
(ic_ic_l)(ic_kc_i)(ic_ic_j) = - (ic_ic_j)(ic_kc_i)(ic_ic_l),
\]
or equivalently
\begin{equation}\label{eq.M-rep_swap}
    \mathcal{M}_{\overleftarrow{il}} \mathcal{M}_{\overleftarrow{kij}}
    =
    -\mathcal{M}_{\overleftarrow{ij}}
    \mathcal{M}_{\overleftarrow{kil}}.
\end{equation}
Comparing \eqref{eq.M-rep_swap} with \eqref{eq.S-rep-swap} shows the identity between the $\mathcal{S}$-representation and the Majorana representation.


